%% kapa: exam=1 round=24 year=2013 period=1 form=A subject=civil_law subject_ko=민법 seq=15 number=15
%% source: Q-Net 감정평가사 기출문제 (cst00309) artlSeq=5073802 file=extracted/1008/5073802/2013년 제24회 감정평가사 1차시험 기출문제/1교시(A형).hwp
%% conversion: kordoc 4.24.51 -> tools/gen1.py (auto); review=proofread-sonnet; flags=none
\begin{question}{15}
표현대리에 관한 설명으로 옳지 않은 것은? (다툼이 있으면 판례에 의함)
\begin{choices}
  \item 표현대리가 성립하는 경우에는 상대방에게 설령 과실이 있더라도 과실상계의 법리를 유추적용할 수 없다.
  \item 등기신청을 위한 대리권을 가진 대리인이 그 권한을 넘어 대물변제를 한 경우 표현대리의 법리는 적용되지 않는다.
  \item 표현대리가 성립한다고 하여 무권대리의 성질이 유권대리로 전환되는 것은 아니다.
  \item 호텔 등의 시설이용 우대회원 모집계약을 체결하면서 자신의 판매점이나 총대리점 등의 명칭을 사용하여 회원모집 안내를 하는 것을 승낙 또는 묵인하였다면 대리권 수여의 표시가 있은 것으로 볼 수 있다.
  \item 대리권 소멸후의 표현대리가 인정되는 경우 그 권한을 넘는 대리행위가 있을 때에는 권한을 넘은 표현대리가 성립할 수 있다.
\end{choices}
\end{question}
\answer{2}
